% Generated draft: requires mathematical and source review.
\documentclass[11pt]{article}
\usepackage[margin=25mm]{geometry}
\usepackage{fontspec}
\setmainfont{Latin Modern Roman}
\usepackage{amsmath,amssymb,mathtools,unicode-math}
\setmathfont{Latin Modern Math}
\usepackage{enumitem,needspace,xurl,hyperref}
\setlength{\parindent}{0pt}
\setlength{\parskip}{5pt}
\setlength{\emergencystretch}{3em}
\newcommand{\N}{\mathbb{N}}
\newcommand{\Z}{\mathbb{Z}}
\newcommand{\Q}{\mathbb{Q}}
\newcommand{\R}{\mathbb{R}}
\newcommand{\C}{\mathbb{C}}
\newcommand{\F}{\mathbb{F}}
\newcommand{\A}{\mathbb{A}}
\newcommand{\PP}{\mathbb{P}}
\newcommand{\E}{\mathbb{E}}
\newcommand{\eps}{\varepsilon}
\newcommand{\dd}{\,\mathrm{d}}
\begin{document}
\section*{UnsolvedMath ID: 700973}
\textbf{Source problem code:} catalogue entry
\subsection{Definitions and mathematical statement}
For groups \(G\) and \(H\), an isomorphism \(f\colon G\to H\) is a bijection satisfying
\[
\forall x,y\in G,\qquad f(xy)=f(x)f(y).
\]
The term \emph{Golod group} is used in the sense referenced by Problem 9.76 of the Kourovka Notebook. The term \(AT\)-\emph{group} is used in the sense of the definition in the cited paper of A. V. Rozhkov. Neither definition is reproduced in the supplied record, so these two conditions remain dependent on the referenced sources; no additional hypotheses or interpretation of either term are imposed here.

The question is whether there exist groups \(G\) and \(H\) and an isomorphism \(f\colon G\to H\) such that \(G\) is a Golod group in the referenced sense and \(H\) is an \(AT\)-group in the referenced sense. A self-contained formulation requires the missing definitions.
\subsection{Short English statement}
Does there exist a Golod group, in the sense referenced by Kourovka Problem 9.76, that is isomorphic to an \(AT\)-group as defined in the cited paper of A. V. Rozhkov? The supplied record does not reproduce either definition.
\subsection{Sources}
\begin{itemize}
\item Local UnsolvedMath registry: supplied complete record with permanent ID \(700973\), exact source code \texttt{catalogue entry}, titled \emph{Kourovka Notebook Problem 13.55}. Its supplied status is \emph{open}; this label is not an independently verified mathematical or literature assessment.
\item \emph{The Kourovka Notebook (Unsolved Problems in Group Theory)}, Problem 13.55, proposed by A. V. Timofeenko in 1995, 13th issue (1995). The record also identifies the 21st-issue edition (2026), \texttt{arXiv:1401.0300v48}, and names E. I. Khukhro and V. D. Mazurov as editors of the collection.
\item Kourovka Notebook Problem 9.76: the supplied statement's cross-reference for the meaning of a Golod group. Its contents are not supplied.
\item A. V. Rozhkov, \emph{Math. Notes}, volume 40, no. 5 (1986), pp. 827--836: the supplied reference for the definition of an \(AT\)-group. The record supplies neither the article title nor the definition.
\end{itemize}
\end{document}
